\subsubsection*{Several-insertion expansion}
For a finite set of slots $I_0$, work in the commutative algebra $\C[z_i:i\in I_0]/(z_i^2:i\in I_0)$. Write $z_I=\prod_{i\in I}z_i$. The exact source density starts at
\begin{equation}\label{eq:E-multiple-density}
 \mathcal H_N(z;U)=\density\rho_N(U)\prod_{i\in I_0}(1+z_iF_i(U)),
 \qquad \mathcal H_m=\blockmap{m}_{\beta_{N,p,m+1}}\mathcal H_{m+1}.
\end{equation}
Thus its normalized $z_I$ integral is the cutoff moment of the corresponding physical insertions. At level~$m$ write its local representation as
\begin{equation}\label{eq:E-source-form}
 e^{\Phi_m(U)+\sum_i z_if_{i,m}(U)+\sum_X H_{m,X}(z;U)}
 \sum_{\mathcal A\ {\rm eligible}}
 \prod_{Y\in\mathcal A}
 \left(K_{m,Y}(U)+\sum_i z_i k_{i,m,Y}(U)+K^{\ge2}_{m,Y}(z;U)\right).
\end{equation}
Here $\Phi_m$ is the unmarked exponent, $f_{i,m}$ and $k_{i,m,Y}$ are the complete one-insertion exponential and forced terms, including their physical spacing factors. The polynomials $H_{m,X}$ and $K^{\ge2}_{m,Y}$ contain only modes $z_I$ with $|I|\ge2$. Their support contains every anchor indexed by a nonzero mode. The regular terms include their health masks, and each forced term retains its eligibility test. Formula~\eqref{eq:E-source-form} distinguishes the exact multiple-source density from the one obtained by keeping only its unmarked and singleton data.

The next estimate controls one insertion in every reflected cell. Its normalization cost is proportional to the number of marks and to the number of coarse sites, rather than to the number of fine sites.
\begin{lemma}[Replicated source bound]\label{lem:E-replication}
For every $d$ and bounded coefficient ball there are radii $\mathfrak r_k>0$, independent of~$p$, such that, for $m\ge r(d)$, $k\ge m$ and $B_k^2\le a_m/a_k$, eventually in~$N>k$, uniformly in~$p$,
\begin{equation*}
 \left|\left\langle\prod_{C\in\mathcal C^{\rm tile}_m}(O(c))_C\right\rangle_{N,p}\right|
 \mathfrak r_k^{|\mathcal C^{\rm tile}_m|}
 \le e^{|\Lambda_{n_{p,k}}|q_k+6|\mathcal C^{\rm tile}_m|}.
\end{equation*}
For every $\eta>0$, eventually $6-\log\mathfrak r_k\le(d+\eta)s_k$.
\end{lemma}
\begin{proof}
Use Definitions~\ref{def:D05a}--\ref{def:D08}, Lemmas~\ref{lem:E1.1}, \ref{lem:E-marked-step} and~\ref{lem:E2.2a}, and the expansion in Lemma~\ref{lem:D-connected}. Enlarge the fixed coefficient radius to include its signed and conjugated images. Put
\begin{equation*}
 E_k=R' C_{J_k}\prod_{i<k}C_{J_{i+1}},\qquad
 \mathfrak r_k=a_k^d\widehat\beta_k^{-w_d}/E_k,
\end{equation*}
where $R'>0$ depends on the coefficient ball and $w_d$ is a fixed nonnegative integer large enough for its finite polynomial evaluations and forced counterparts. The transport inequality of Lemma~\ref{lem:E2.2}, iterated from $r(d)$, bounds the level-$k$ coefficients by the product in $E_k$. Threshold powers introduce only a fixed power of $\widehat\beta_k$; enlarging $w_d$ includes them. Hence multiplying every mark by $\mathfrak r_k$ bounds its full singleton size by the positive source allowance $\widehat\beta_k^{-w_d}\le1$.

For the local expansion, give a mode $z_I$ the weight $\prod_{i\in I}w_i$, with $w_i>0$. The resulting coefficient norm is submultiplicative: multiplication sums over disjoint splittings of~$I$, exactly the splittings in the product of the two weighted norms. The expansion of Lemma~\ref{lem:D-connected} therefore applies coefficientwise. Its connected terms with at least two marks have the additional support estimate
\begin{equation}\label{eq:E-marked-cover}
 B_k|I|\le64\cover_{B_k}(X),\qquad |I|\ge2,
\end{equation}
when distinct anchors are separated by $B_k^2$. To verify it, take a minimal cover by $M$ radius-$B_k$ balls with connected union. Join their centers by a tree with edges of length at most $2B_k+1$. Its doubled traversal has length at most $6B_k(M-1)$. Assign each anchor to one covering center. Distinct anchors have distinct centers, and the distance between any two assigned centers is at least $B_k^2-2B_k$. Visit these centers in their cyclic order in the traversal. The triangle inequality gives $|I|(B_k^2-2B_k)\le6B_k(M-1)$. For $B_k\ge8$ this implies the stated bound, with room in the constant. The argument uses the periodic max-distance and needs no nonwrapping assumption. This is the marked covering estimate used in the support norm of Lemma~\ref{lem:D-connected}.

It permits each source weight to be enlarged by $e^{\delta B_k/64}$ when $\delta$ units of the support exponent are removed. Singleton terms are bounded directly by their one-mark norms. Taylor's product rule introduces at most the factorial of the derivative order; the stage column constant was chosen above this factorial. With $j=J_{k+1}$, Lemma~\ref{lem:D-connected} supplies the rooted derivative bound at support exponent $\sigma_{J_k}+3$ and the forced bound $C_j e^{-r_j\beta_{N,p,k+1}t_{k+1}^2}$ at exponent $\sigma_{J_k}+2$. Multiplying by a mode weight changes only the marked letter in its tree sum. Item~\ref{it:E-columns} pays the derivative factorial. The weighted sums therefore retain those bounds; decreasing the final exponent from $100$ to any value in $[98,100]$ only decreases the sums. The forced normalization uses half the suppression, leaving the other half for the marked derivative and attachment losses. Their spatial anchor is always one of the marks; a nonzero mode is assigned only to a support containing all its marked anchors. This proves the support assertion in~\eqref{eq:E-source-form}, as well as its absolute convergence. Restricting the set of marks commutes with the construction, since every operation takes disjoint unions of their modes.

The reflected anchors at level~$k$ are the embedded sites of the level-$m$ torus. They have spacing $a_m/a_k\ge B_k^2$. Their parents agree under each subsequent block map. The inequality persists at all finer levels because $B_{i+1}^2\le2B_i^2$ and $\ell_i\ge3$. Occupancy is therefore one at every such level. These are precisely the conditions needed for the preceding weighted sums; the number of marks does not enter their local bounds.

At the stopping level, the weighted supremum of each complete singleton exponential term is at most $3$, and that of its forced profile is at most~$1$. Indeed, its polynomial part is at most its coefficient bound, its marked regular sum is at most that bound, and the normalization of the forced part is at most one; Lemma~\ref{lem:E2.2} and the definition of $E_k$ bound the sum by~$3$. The forced profile follows from its $\kappa_k$ normalization. For~\eqref{eq:E-source-form}, the local absolute sums give an additional two per mark: exponentiation of the connected regular terms and the eligible forced sum each contribute at most one. Integrating the remaining scalar-free unmarked exponent and dividing by the unmarked partition function gives at most $e^{|\Lambda_{n_{p,k}}|q_k}$, with $q_k$ as in~\eqref{eq:E-vacuum-majorant}. Thus coefficient extraction in the weighted source norm yields
\begin{equation*}
 \mathfrak r_k^{|I_0|}\left|\left\langle\prod_{i\in I_0}F_i\right\rangle_{N,p}\right|
 \le\exp\{|\Lambda_{n_{p,k}}|q_k+(3+1+2)|I_0|\}.
\end{equation*}
Exact integration preserves these normalized coefficients, proving the first assertion. The ratio of site and cell counts is $(a_m/a_k)^4$.

Finally~\eqref{eq:E-log-cost}, summed as in~\eqref{eq:E-running}, gives
\begin{equation*}
 \sum_{i<k}\log C_{J_{i+1}}=o(s_k),\qquad
 \log\widehat\beta_k=o(s_k).
\end{equation*}
The separate final column is included in this sum for $k>0$. Hence, for every $\eta>0$, eventually $6-\log\mathfrak r_k\le(d+\eta)s_k$. All choices precede the period, and all the estimates are on full cutoff tails.
\end{proof}
